\section{Proofs}
\label{sec:supp_proof}

\subsection{Proof of the allocation stability result}
\label{sec:supp_proof_stability}
This is the proof of Theorem~1 of the main paper (Allocation stability), repeated here in full.
Write $s_k = -1$ for $k \in \{1,5\}$ and $s_k = +1$ otherwise, so that the allocation equation of the
main paper reads $a(\mathbf{r}) = b + \sum_{k} s_k C_k w_k r_k$ with $b$ independent of $\mathbf{r}$,
where $\mathbf{r} = (r_1, \ldots, r_5) \in [0,1]^5$ is the vector of rule truth values at one pixel,
$C_k$ are the fixed budgets, and $w_k = \sigma(\theta_k)$ are the learned rule weights.

\begin{proof}
Subtracting two such expressions gives
\begin{equation*}
a(\mathbf{r}) - a(\mathbf{r}') = \sum_{k} s_k C_k w_k (r_k - r'_k).
\end{equation*}
Taking absolute values and applying the triangle inequality,
\begin{equation*}
\begin{split}
\big| a(\mathbf{r}) - a(\mathbf{r}') \big|
&\le \sum_{k} C_k w_k \, |r_k - r'_k| \\
&\le \Big( \sum_{k} C_k w_k \Big) \lVert \mathbf{r} - \mathbf{r}' \rVert_\infty ,
\end{split}
\end{equation*}
which is the stability bound of Definition~1 of the main paper with $L = \sum_k C_k w_k$ as claimed.
The bound is attained by taking $r_k - r'_k = s_k t$ for any $t$ small enough to keep both vectors in
the unit cube, since then every term of the sum has the same sign. Because each $C_k$ is positive and
each $w_k = \sigma(\theta_k)$ lies strictly in $(0,1)$, the sum is strictly less than
$\sum_k C_k = 0.75$.
\end{proof}

\subsection{Proof of the contrast margin result}
\label{sec:supp_proof_cnr}
This is the proof of Theorem~2 of the main paper, repeated here in full. Write
$\mu_\mathcal{T}$ and $\mu_\mathcal{B}$ for the mean over tissue and over background, and
$\sigma_\mathcal{B}$ for the standard deviation over background. Let $\mathbf{b}$ be the backbone
output, $\mathbf{q}$ the candidate, and $\hat{\mathbf{x}} = (1-w)\mathbf{b} + w\mathbf{q}$ the output
with $w \in (0,1]$.

\begin{theorem}[Contrast margin]
Let $\eta_\mathcal{T} = \mu_\mathcal{T}(\mathbf{q}) - \mu_\mathcal{T}(\mathbf{b}) \ge 0$ and let
$\eta_\mathcal{B} \ge \mu_\mathcal{B}(\mathbf{q}) - \mu_\mathcal{B}(\mathbf{b})$. If
$\sigma_\mathcal{B}(\mathbf{q}) \le \sigma_\mathcal{B}(\mathbf{b})$, that
$\sigma_\mathcal{B}(\mathbf{b}) > 0$, that the tissue and background sets are the same for all three
images, and that the bounded numerator is non negative, then
\begin{equation}
\mathrm{CNR}(\hat{\mathbf{x}}) \ge
\frac{\big(\mu_\mathcal{T}(\mathbf{b}) - \mu_\mathcal{B}(\mathbf{b})\big)
      + w\,(\eta_\mathcal{T} - \eta_\mathcal{B})}{\sigma_\mathcal{B}(\mathbf{b})} .
\end{equation}
\end{theorem}

\begin{proof}
The mean is linear, so over any fixed set of pixels the mean of a convex combination is the same
convex combination of the means. Applying this over the tissue set and over the background set,
\begin{align*}
\mu_\mathcal{T}(\hat{\mathbf{x}}) &= (1-w)\,\mu_\mathcal{T}(\mathbf{b}) + w\,\mu_\mathcal{T}(\mathbf{q})
 = \mu_\mathcal{T}(\mathbf{b}) + w\,\eta_\mathcal{T}, \\
\mu_\mathcal{B}(\hat{\mathbf{x}}) &= \mu_\mathcal{B}(\mathbf{b}) + w\,\eta_\mathcal{B}',
\end{align*}
where $\eta_\mathcal{B}' = \mu_\mathcal{B}(\mathbf{q}) - \mu_\mathcal{B}(\mathbf{b}) \le \eta_\mathcal{B}$
by hypothesis. Subtracting,
\begin{equation*}
\begin{split}
\mu_\mathcal{T}(\hat{\mathbf{x}}) - \mu_\mathcal{B}(\hat{\mathbf{x}})
&= \big(\mu_\mathcal{T}(\mathbf{b}) - \mu_\mathcal{B}(\mathbf{b})\big) + w\,(\eta_\mathcal{T} - \eta_\mathcal{B}') \\
&\ge \big(\mu_\mathcal{T}(\mathbf{b}) - \mu_\mathcal{B}(\mathbf{b})\big) + w\,(\eta_\mathcal{T} - \eta_\mathcal{B}),
\end{split}
\end{equation*}
which lower bounds the numerator.

For the denominator, every background pixel of $\hat{\mathbf{x}}$ is the same convex combination of
the corresponding pixels of $\mathbf{b}$ and $\mathbf{q}$. Treating the background pixel values as a
random variable and applying Minkowski's inequality to the centered variables,
\begin{equation*}
\begin{split}
\sigma_\mathcal{B}(\hat{\mathbf{x}}) &\le (1-w)\,\sigma_\mathcal{B}(\mathbf{b}) + w\,\sigma_\mathcal{B}(\mathbf{q}) \\
&\le (1-w)\,\sigma_\mathcal{B}(\mathbf{b}) + w\,\sigma_\mathcal{B}(\mathbf{b}) = \sigma_\mathcal{B}(\mathbf{b}),
\end{split}
\end{equation*}
using the hypothesis $\sigma_\mathcal{B}(\mathbf{q}) \le \sigma_\mathcal{B}(\mathbf{b})$ in the second
step. Dividing the lower bound on the numerator by the upper bound on the denominator gives the
claim. Since $w > 0$ in every branch of the blending rule, the bound is at least the baseline
contrast whenever $\eta_\mathcal{T} > \eta_\mathcal{B}$.
\end{proof}

\noindent The results section of the main paper reports how often each of the three hypotheses holds on the
test set, so the reader can see when the guarantee is active rather than assuming it always is.
